\chapter{Measuring Reactivity and Feedback Coefficients}
\label{ch:measure}

A reactor's stability rests on the signs and magnitudes of its feedback
coefficients. These are not merely calculated; they are \emph{measured}, both
during commissioning and periodically through life. This chapter surveys the
experimental methods, which double as a practical illustration of the kinetics
theory and underline a sobering point: the RBMK's dangerous coefficients were, in
principle, measurable and partly known.

\section{Static and dynamic reactivity measurement}

\textbf{Rod-drop and source-jerk methods} infer reactivity from the transient
following a sudden change. In a rod-drop, a calibrated rod is dropped and the
prompt power decrease is fitted to the point-kinetics solution to extract the
inserted negative reactivity. \textbf{Rod-oscillator methods} insert a sinusoidal
reactivity and measure the power response; the amplitude and phase give the
zero-power transfer function $Z(j\omega)$ of Chapter~\ref{ch:linear} directly,
mapping the reactor's dynamic response across frequency. \textbf{Inverse
kinetics}---solving the point-kinetics equations ``backwards'' for $\rho(t)$ given a
measured power trace---is the basis of the modern \emph{reactivity meter}, a
real-time instrument used during physics testing.

\section{Period measurement and the inhour calibration}

The simplest reactivity measurement is the \emph{positive-period} method: insert a
small known reactivity, measure the resulting stable period, and confirm it against
the inhour equation \eqref{eq:inhour}. This both measures reactivity and validates
the kinetics data. The agreement between measured periods and the inhour prediction
(Figure~\ref{fig:inhour}) is, in effect, a routine check that the reactor's
delayed-neutron parameters and generation time are as expected.

\section{Measuring temperature and void coefficients}

The \textbf{isothermal temperature coefficient} is measured at low power by slowly
heating the whole coolant/moderator system (using pump heat, with the reactor just
critical) and recording the compensating rod motion: the rod worth needed to hold
criticality per degree is the coefficient. The \textbf{moderator temperature
coefficient (MTC)} and \textbf{Doppler (fuel) coefficient} can be separated by
exploiting their different timescales---the fuel responds promptly, the moderator
with the coolant transport delay---using power-oscillation or noise techniques. The
\textbf{void coefficient} is inferred from controlled changes in boiling or
moderator level, or in test reactors by deliberately introducing voids. The
\textbf{power coefficient} is measured directly by recording the reactivity
(via the reactivity meter) required to change power.

\section{Noise analysis: listening to the reactor}

A powerful, non-invasive technique is \emph{reactor noise analysis}. The neutron
flux fluctuates around its mean due to the statistical nature of fission and the
reactor's dynamic response; the spectrum of these fluctuations encodes the
transfer function and hence the feedback. Measuring the \emph{decay ratio}
(Chapter~\ref{ch:linear}) from the autocorrelation of small flux oscillations is
the standard method for monitoring BWR stability margins online---a direct,
continuous read on how close the reactor is to the oscillatory stability boundary.

\begin{remark}[The knowable danger]
The RBMK's positive void coefficient and the positive scram effect were, in
principle, measurable quantities, and the scram anomaly had in fact been observed
at the Ignalina plant in 1983. The tragedy is not that the danger was unmeasurable
but that the measurements and observations were not translated into design fixes,
clear operating limits, and operator understanding. Measurement is only the first
step; acting on it is safety.
\end{remark}

\keypoint{Feedback coefficients are measured by rod-drop, oscillator, inverse
kinetics, isothermal heating, and noise methods, and validated against the inhour
equation. The RBMK's destabilising coefficients were measurable and partly
observed; the failure was organisational---not converting knowledge into design
changes and enforced limits.}

\section*{Exercises}
\addcontentsline{toc}{section}{Exercises}
\begin{enumerate}[leftmargin=*,itemsep=2pt]
\item Describe the rod-drop method for measuring reactivity.
\item How does a rod oscillator measure the zero-power transfer function?
\item Outline how the isothermal temperature coefficient is measured.
\item Explain how reactor noise analysis monitors BWR stability margins.
\item Discuss why the RBMK danger was knowable yet not acted upon.
\end{enumerate}
